\section{Markings, images and kernels}\label{markings-images-and-kernels}

A marked group is a group \(G\) together with a finite indexed
generating family \((r_i)_{i\in I}\). Except where stated, the markers
need not be involutions or distinct. For \(S\subseteq I\), put \[
G_S=\langle r_i:i\in S\rangle,\qquad G_\varnothing=1.
\] We call these subgroups parabolics, retaining their index sets even
when two of them happen to coincide. The marking has the intersection
property, or IP, if \[
G_S\cap G_T=G_{S\cap T}\quad(S,T\subseteq I).
\tag{M1}
\] A marked homomorphism sends each source marker to the corresponding
target marker. Its image on \(G_S\) is therefore the target parabolic
with index set \(S\).

All subgroup products below are products of sets unless a subgroup
assertion is made. We write \(\gamma_1(G)=G\) and
\(\gamma_{j+1}(G)=[\gamma_j(G),G]\). We use the elementary
correspondence and isomorphism theorems for groups. In the finite
\(p\)-group arguments we retain Burnside's basis theorem, the nilpotence
of finite \(p\)-groups, and the usual multilinear commutator description
of the lower-central graded Lie ring. Precisely, its degree-one part
generates the graded Lie ring, and degree \(j\) is generated by
left-normed commutators of weight \(j\) in any generating set. In
the amalgam section we explicitly retain the normal-form theorem for an
amalgamated free product. None of these classical inputs is presented as
a result originating in this manuscript.

\textbf{Lemma UA.K0 (lifting one intersection).} Let \(q:H\to Q\) be a
homomorphism and let \(C\leq A\cap B\), where \(A,B,C\leq H\). If \[
q(A)\cap q(B)=q(C),
\] then \[
A\cap B=C\quad\Longleftrightarrow\quad A\cap B\cap\ker q\leq C.
\tag{M2}
\]

\emph{Proof.} The forward implication is immediate. For the converse,
let \(x\in A\cap B\). Its image belongs to \(q(C)\), so choose
\(c\in C\) with \(q(c)=q(x)\). Then \(xc^{-1}\in A\cap B\cap\ker q\),
which is contained in \(C\) by hypothesis. Multiplying by \(c\) gives
\(x\in C\). The inclusion \(C\leq A\cap B\) gives equality. \(\square\)

\textbf{Theorem UA.K (kernel profiles).} Let \(q:H\twoheadrightarrow Q\)
be a marked epimorphism and suppose that the marking of \(Q\) has IP.
Write \(K=\ker q\) and \(K_S=K\cap H_S\). Then \(H\) has IP if and only
if \[
K_S\cap K_T=K_{S\cap T}\quad(S,T\subseteq I).
\tag{M3}
\]

\emph{Proof.} Apply UA.K0 to \(A=H_S\), \(B=H_T\), and
\(C=H_{S\cap T}\). The image hypothesis is exactly IP in \(Q\). The
subgroup on the left of the containment in (M2) is \(K_S\cap K_T\). Its
containment in \(H_{S\cap T}\) is equivalent to its containment in
\(K_{S\cap T}\); the reverse inclusion always holds. This proves each
required intersection and its converse. \(\square\)

The next result concerns the opposite direction. Passing to a quotient
can create new intersections, and this defect is a product condition in
its kernel.

\textbf{Theorem UA.Q (exact intersections of images).} Let \(q:H\to Q\)
be a homomorphism with kernel \(N\), and let \(A,B\leq H\). Then \[
q(A)\cap q(B)=q(A\cap B)
\quad\Longleftrightarrow\quad
N\cap AB=(N\cap A)(N\cap B).
\tag{M4}
\] If \(q\) is injective on \(A\) and \(B\), the right side reduces to
\(N\cap AB=\{1\}\). If \(N\leq A\), the equality always holds.

\emph{Proof.} Suppose the image intersection is exact, and take
\(ab\in N\), where \(a\in A\) and \(b\in B\). Since \(q(a)=q(b)^{-1}\),
choose \(c\in A\cap B\) with this image. Then \[
ab=(ac^{-1})(cb)\in(N\cap A)(N\cap B).
\] The reverse inclusion requires no hypothesis. Conversely, if
\(q(a)=q(b)\), then \(ab^{-1}\in N\cap AB\). Write \(ab^{-1}=uv\), with
\(u\in N\cap A\) and \(v\in N\cap B\). The element \(u^{-1}a=vb\) lies
in \(A\cap B\) and has image \(q(a)\), proving the required image
equality. Under injectivity both factors on the right of (M4) are
trivial. If \(N\leq A\), both sides of that equality are \(N\).
\(\square\)
